\documentclass[12pt, a4paper]{article}

% ── PACKAGES ──────────────────────────────────────────────────
\usepackage[utf8]{inputenc}
\usepackage[T1]{fontenc}
\usepackage{geometry}
\geometry{margin=2.5cm}
\usepackage{hyperref}
\hypersetup{
    colorlinks=true,
    linkcolor=blue,
    urlcolor=blue,
    citecolor=blue,
    pdftitle={Registered Amendment A1 to ATTRACTOR-LIFE-PREREGISTRATION-v1.0:
              Bilateral Symmetry as a Two-Axis Attractor and
              Revised Dickinsonia Protocol},
    pdfauthor={Eric Robert Lawson},
    pdfkeywords={attractor geometry, LECA, Dickinsonia, Ediacaran,
                 Wnt, BMP, bilateral symmetry, two-axis attractor,
                 dimensional expansion, Reproduction Theorem,
                 registered amendment, OrganismCore}
}
\usepackage{amsmath}
\usepackage{amssymb}
\usepackage{amsthm}
\usepackage{booktabs}
\usepackage{longtable}
\usepackage{array}
\usepackage{parskip}
\usepackage{microtype}
\usepackage{authblk}
\usepackage{titlesec}
\usepackage{fancyhdr}
\usepackage{xcolor}
\usepackage{float}
\usepackage{graphicx}
\usepackage{setspace}
\usepackage{enumitem}
\usepackage{tabularx}
\usepackage{multirow}
\usepackage{doi}

% ── COLOURS ───────────────────────────────────────────────────
\definecolor{repobox}{RGB}{240,247,255}
\definecolor{lockbox}{RGB}{245,255,245}
\definecolor{warnbox}{RGB}{255,248,220}
\definecolor{amendbox}{RGB}{255,240,240}
\definecolor{theorembox}{RGB}{248,240,255}

% ── THEOREM ENVIRONMENTS ──────────────────────────────────────
\newtheorem{theorem}{Theorem}[section]
\newtheorem{corollary}[theorem]{Corollary}
\newtheorem{lemma}[theorem]{Lemma}
\theoremstyle{definition}
\newtheorem{definition}{Definition}[section]
\newtheorem{prediction}{Prediction}[section]
\theoremstyle{remark}
\newtheorem{remark}{Remark}[section]

% ── HEADER / FOOTER ───────────────────────────────────────────
\pagestyle{fancy}
\fancyhf{}
\rhead{\small OrganismCore \textbar{} 2026-03-12}
\lhead{\small Amendment A1 --- Dickinsonia Protocol Revision}
\rfoot{\small Page \thepage}
\renewcommand{\headrulewidth}{0.4pt}

% ── SECTION FORMATTING ────────────────────────────────────────
\titleformat{\section}
  {\large\bfseries}{\thesection}{1em}{}[\titlerule]
\titleformat{\subsection}
  {\normalsize\bfseries}{\thesubsection}{1em}{}

% ══════════════════════════════════════════════════════════════
% TITLE
% ═══════════════════════════════���══════════════════════════════
\title{
    \vspace{-1em}
    \textbf{Registered Amendment A1}\\[0.5em]
    \large\textbf{to ATTRACTOR-LIFE-PREREGISTRATION-v1.0}\\[0.6em]
    \normalsize\textit{Bilateral Symmetry as a Two-Axis Attractor:}\\
    \normalsize\textit{Geometric Derivation, Literature Confirmation,}\\
    \normalsize\textit{and Revised Dickinsonia Protocol}\\[0.6em]
    \small OrganismCore Document:
    ATTRACTOR-LIFE-PREREGISTRATION-v1.0-AMENDMENT-A1
}

\author[1]{Eric Robert Lawson}
\affil[1]{
    OrganismCore\\
    \href{mailto:OrganismCore@proton.me}
         {OrganismCore@proton.me}\\
    ORCID:
    \href{https://orcid.org/0009-0002-0414-6544}
         {0009-0002-0414-6544}\\
    Repository:
    \href{https://github.com/Eric-Robert-Lawson/attractor-oncology}
         {github.com/Eric-Robert-Lawson/attractor-oncology}
}

\date{
    Amendment date: 2026-03-12\\
    \small\textit{This amendment is filed before any
    experimental confirmation. It constitutes a registered
    prediction refinement, not a framework revision.}
}

% ══════════════════════════════════════════════════════════════
\begin{document}
% ══════════════════════════════════════════════════════════════

\maketitle
\thispagestyle{fancy}

% ── AMENDMENT IDENTITY BOX ────────────────────────────────────
\begin{center}
\colorbox{amendbox}{
\begin{minipage}{0.88\textwidth}
\vspace{0.6em}
\centering
\textbf{REGISTERED AMENDMENT A1}\\[0.4em]
\textbf{Amends:}
ATTRACTOR-LIFE-PREREGISTRATION-v1.0\\
\textbf{Original pre-registration DOI:}
\href{https://doi.org/10.5281/zenodo.18986790}
     {\texttt{https://doi.org/10.5281/zenodo.18986790}}\\[0.3em]
\textbf{Companion publication DOI (LUCA paper):}
\href{https://doi.org/10.5281/zenodo.18988844}
     {\texttt{https://doi.org/10.5281/zenodo.18988844}}\\[0.3em]
\textbf{This amendment DOI:} [assigned by Zenodo on upload]\\[0.3em]
\small
Amendment type: \textbf{Prediction refinement}\\
Framework status: \textbf{Unchanged --- no framework revision}\\
Original pre-registration: \textbf{Permanently locked ---
not modified}\\
Experiment status at time of amendment:
\textbf{Not yet run}\\[0.3em]
\textit{The original pre-registration record is not amended.
It stands as published. This document refines one prediction
within it (the Dickinsonia Protocol) using new geometric
information derived by integrating the Wnt/BMP literature
into the attractor framework. All other predictions in
the original are unchanged.}
\vspace{0.5em}
\end{minipage}
}
\end{center}

\vspace{0.8em}

% ── ABSTRACT ──────────────────────────────────────────────────
\begin{abstract}
\noindent
This document constitutes Registered Amendment A1 to the
pre-registration record ATTRACTOR-LIFE-PREREGISTRATION-v1.0
(Zenodo DOI: \href{https://doi.org/10.5281/zenodo.18986790}
{10.5281/zenodo.18986790}).

The amendment is triggered by the geometric integration of
two bodies of literature --- the Wnt signalling evolution
review (\textit{Nature Reviews Genetics}, 2024) and the
BMP/Wnt orthogonal axis integration papers (PLOS Genetics,
2023; \textit{Discover Developmental Biology}, 2024) ---
into the Attractor Geometry framework. This integration
reveals that the original Dickinsonia Protocol's prediction
of bilateral Ediacaran morphology from Wnt activation alone
is under-specified.

The geometric derivation produces the following result: bilateral
symmetry is not a single-axis attractor. It is a
\textbf{two-axis attractor} constituted by the orthogonal
intersection of the Wnt/$\beta$-catenin axis
(anterior-posterior, Axis~1) and the BMP/ADMP axis
(dorsal-ventral, Axis~2). The transition from pre-bilateral
Ediacaran morphology to bilateral Ediacaran morphology is a
\textbf{dimensional expansion event} in the body plan
attractor space --- structurally identical to the
mitochondrial dimensional expansion (the FECA$\to$LECA
event) but operating in morphological rather than metabolic
attractor space.

This derivation produces two consequences: (1)~the
Dickinsonia Protocol requires two sequential interventions,
not one, to access the bilateral Ediacaran attractor; (2)~the
pre-bilateral Ediacaran attractor (Axis~1 only) is itself
an independently accessible and falsifiable intermediate
stage.

The \textbf{framework is unchanged}. The Reproduction
Theorem, the Inverse Accessibility Theorem, all LECA
predictions, and all other Dickinsonia predictions in the
original pre-registration stand without modification. This
amendment refines the protocol and introduces two new
independently falsifiable predictions derived from the
revised geometry.

Amendment filed before any experimental confirmation.
Timestamp: 2026-03-12.
\end{abstract}

\vspace{0.5em}
\noindent\textbf{Keywords:}
registered amendment; attractor geometry; Dickinsonia;
bilateral symmetry; two-axis attractor; Wnt signalling;
BMP signalling; dimensional expansion; Ediacaran;
Reproduction Theorem; prediction refinement; OrganismCore

\newpage
\tableofcontents
\newpage

% ══════════════════════════════════════════════════════════════
\section{Document Metadata}
% ═══════════════════���═════════════════════════════���════════════

\begin{center}
\begin{tabular}{ll}
\toprule
\textbf{Field} & \textbf{Value} \\
\midrule
Document ID
  & ATTRACTOR-LIFE-PREREGISTRATION-v1.0-AMENDMENT-A1 \\
Amendment type
  & Prediction refinement (protocol specification) \\
Amends
  & ATTRACTOR-LIFE-PREREGISTRATION-v1.0 \\
Original DOI
  & \href{https://doi.org/10.5281/zenodo.18986790}
         {10.5281/zenodo.18986790} \\
Companion DOI (LUCA paper)
  & \href{https://doi.org/10.5281/zenodo.18988844}
         {10.5281/zenodo.18988844} \\
Date
  & 2026-03-12 \\
Author
  & Eric Robert Lawson \\
Organisation
  & OrganismCore \\
ORCID
  & \href{https://orcid.org/0009-0002-0414-6544}
         {0009-0002-0414-6544} \\
Contact
  & \href{mailto:OrganismCore@proton.me}
         {OrganismCore@proton.me} \\
Framework status
  & Unchanged \\
Experiment status
  & Not yet run at time of filing \\
Biological contradictions (cumulative)
  & 0 \\
Type B failures (cumulative)
  & 0 \\
\bottomrule
\end{tabular}
\end{center}

% ══════════════════════════════════════════════════════════════
\section{Amendment Scope and Epistemic Status}
% ══════════════════════════════════════════════════════════════

\subsection{What this amendment does}

This amendment does precisely one thing:

\textbf{It refines the Dickinsonia Protocol} in the original
pre-registration by specifying the two-step intervention
required to access the bilateral Ediacaran attractor, and
by introducing the pre-bilateral Ediacaran attractor as an
independently accessible and falsifiable intermediate.

\subsection{What this amendment does not do}

\begin{itemize}
  \item It does not modify the Reproduction Theorem.
  \item It does not modify the Inverse Accessibility Theorem.
  \item It does not modify any LECA prediction.
  \item It does not modify the Cross-Kingdom Convergence
        Prediction.
  \item It does not modify the LINE-1 independent access
        prediction.
  \item It does not modify the LUCA geometric status.
  \item It does not modify the falsifiability conditions F1--F7.
  \item It does not touch the original pre-registration
        document. That document is permanently locked and
        stands as published.
\end{itemize}

\subsection{Framework vs.\ prediction: the critical
distinction}

A \textbf{framework invalidation} would require that the
geometry of developmental attractors is incorrect --- that
the Reproduction Theorem is false, that developmental stages
are not compressed evolutionary attractors, or that
arrest does not expose ancestral morphology.

A \textbf{prediction refinement} occurs when the framework
is correct but a prediction derived from it was
under-specified: the correct attractor was identified, but
the intervention required to reach it was incompletely
described.

This amendment is a prediction refinement.

The original Dickinsonia Protocol correctly identified the
target attractor: bilateral Ediacaran-grade multicellular
morphology. The protocol incorrectly specified that
Wnt/$\beta$-catenin activation alone is the sufficient
intervention to reach it.

The geometry, once the Wnt/BMP literature is integrated,
reveals that bilateral symmetry requires two orthogonal axes
and therefore two sequential interventions. The attractor
was correctly identified. The coordinate system of the
attractor was under-specified.

This is the same logical structure as Mendeleev correctly
identifying that an element with specific properties must
exist, and then refining the atomic weight estimate as
more data became available. The prediction is not wrong.
The specification is more precise.

\subsection{Epistemic significance}

This amendment is filed before any experiment. The
two-step protocol is derived entirely from geometry. No
empirical result has been observed. The amendment is
pre-experimental throughout.

This means: when the experiment runs and the two-step
protocol produces bilateral Ediacaran-grade morphology
(Prediction A1-P2, below), the prediction will have been
filed in full before the result was obtained. The
pre-experimental record is intact and is strengthened
by this amendment, not weakened.

% ══════════════════════════════════════════════════════════════
\section{The Literature That Triggered the Amendment}
% ══════════════════════════════════════════════════════════════

The following literature, integrated into the attractor
framework, produced the geometric derivation in Section~4.

\begin{center}
\begin{tabular}{p{4.5cm}p{7cm}}
\toprule
\textbf{Source} & \textbf{Finding integrated} \\
\midrule
Holzem, Boutros, Holstein.
\textit{Nature Reviews Genetics} (2024).
DOI: 10.1038/s41576-024-00699-w
  & Wnt signalling is the universal mechanism of
    primary axis (AP) specification across all Metazoa.
    Wnt pathway components are present in all animals
    including Placozoa. Wnt specifies Axis~1 only. \\
\addlinespace
Clark \& Petersen.
\textit{PLOS Genetics} (2023).
DOI: 10.1371/journal.pgen.1010608
  & BMP suppresses Wnt to integrate patterning of
    orthogonal body axes. BMP and Wnt are antagonistic
    regulators of two orthogonal axes, both required
    for bilateral symmetry. \\
\addlinespace
M\"{o}rsdorf et al.
\textit{Discover Developmental Biology} (2024).
DOI: 10.1007/s00427-024-00714-4
  & BMP signalling is highly conserved AND highly
    evolvable in secondary axis patterning across
    Cnidaria and Bilateria. BMP pathway components
    predate bilateral symmetry. \\
\addlinespace
Uveira et al.
\textit{eLife} (2025).
DOI: elife.104508
  & In \textit{Clytia hemisphaerica} (cnidarian),
    Wnt3 specifies the primary oral-aboral axis AND
    globally orients planar cell polarity. A bilateral
    axis is NOT produced by Wnt alone: Wnt produces
    Axis~1; bilateral symmetry requires Axis~2. \\
\addlinespace
Maddox \& Azuma (2016); multiple sources (2022--2024)
  & Wnt activation in non-bilaterian organisms
    (Hydra, Trichoplax, cnidarians) produces
    ectopic primary axes, not bilateral symmetry.
    Bilateral symmetry requires the deployment of
    BMP orthogonally to the established Wnt axis. \\
\bottomrule
\end{tabular}
\end{center}

% ══════════════════════════════════════════════════════════════
\section{The Geometric Derivation}
% ══════════════════════════════════════════════════════════════

\subsection{Bilateral symmetry is a two-axis attractor}

\begin{definition}[Single-Axis Attractor]
A single-axis attractor is a stable multicellular body plan
configuration in which one spatial polarity axis is
specified. Organisms in a single-axis attractor have a
defined oral-aboral or anterior-posterior axis but no
orthogonal dorsal-ventral specification. Examples:
\textit{Trichoplax adhaerens} (oral-aboral polarity);
\textit{Hydra} (oral-aboral axis, radial symmetry);
pre-bilateral Ediacaran organisms.
\end{definition}

\begin{definition}[Two-Axis Attractor]
A two-axis attractor is a stable multicellular body plan
configuration in which two orthogonal spatial polarity axes
are simultaneously specified. Organisms in a two-axis
attractor have a defined anterior-posterior axis (Axis~1,
Wnt-mediated) AND a defined dorsal-ventral axis (Axis~2,
BMP-mediated), orthogonal to Axis~1. The intersection of
these two axes defines bilateral symmetry. All Bilateria
occupy two-axis attractors.
\end{definition}

\begin{theorem}[Bilateral Symmetry as Dimensional Expansion]
\label{thm:bilateral}
The evolutionary transition from the pre-bilateral
Ediacaran attractor to the bilateral Ediacaran attractor
is a dimensional expansion of the body plan attractor
space.

\begin{proof}
Let $\mathcal{A}_{E1}$ denote the pre-bilateral Ediacaran
attractor (single-axis: Wnt-mediated oral-aboral or AP
axis; no DV axis). Let $\mathcal{A}_{E2}$ denote the
bilateral Ediacaran attractor (two-axis: Wnt-mediated
Axis~1 plus BMP-mediated Axis~2 orthogonal to Axis~1).

$\mathcal{A}_{E1}$ occupies a one-dimensional spatial
specification space: one axis, one polarity gradient.

$\mathcal{A}_{E2}$ occupies a two-dimensional spatial
specification space: two orthogonal axes, two independent
polarity gradients.

$\mathcal{A}_{E2}$ is not reachable from $\mathcal{A}_{E1}$
by modification of Axis~1 alone. It requires the addition of
a second independent axis. The addition of an independent
dimension to the attractor space is a dimensional expansion
by definition.

Therefore the $\mathcal{A}_{E1} \to \mathcal{A}_{E2}$
transition is a dimensional expansion. $\square$
\end{proof}
\end{theorem}

\begin{remark}
This dimensional expansion is structurally identical to
the FECA$\to$LECA mitochondrial expansion event, which
added an independent energy dimension to the cellular
attractor space. Both events:
\begin{enumerate}[label=(\roman*)]
  \item add a new independent dimension to an existing
        attractor space;
  \item require a specific molecular event to open the
        new dimension (mitochondrial endosymbiosis for
        FECA$\to$LECA; BMP deployment orthogonal to Wnt
        for $\mathcal{A}_{E1} \to \mathcal{A}_{E2}$);
  \item are irreversible under normal conditions once
        the new dimension is stabilised;
  \item produce descendants that are permanently in
        the expanded attractor space.
\end{enumerate}
The geometry is the same. The domain differs.
\end{remark}

\subsection{The Ediacaran attractor has internal structure}

The original pre-registration treated the Ediacaran
attractor as a single target: the bilateral Ediacaran-grade
morphology of \textit{Dickinsonia}.

The geometric derivation reveals that the Ediacaran
attractor has \textbf{two stages}, not one:

\begin{center}
\begin{tabular}{p{3cm}p{4.5cm}p{4.5cm}}
\toprule
\textbf{Stage} & \textbf{$\mathcal{A}_{E1}$:
Pre-bilateral Ediacaran} &
\textbf{$\mathcal{A}_{E2}$: Bilateral Ediacaran
(\textit{Dickinsonia}-grade)} \\
\midrule
Axes & 1 (Wnt: oral-aboral or AP only) &
2 (Wnt: AP + BMP: DV, orthogonal) \\
Symmetry & Radial or disc & Bilateral \\
Body plan & Sheet or disc morphology &
Elongated bilateral form \\
Intervention &
Wnt/$\beta$-catenin activation (Step~1 only) &
Wnt activation (Step~1) \textbf{+}
BMP activation orthogonal to Wnt axis (Step~2) \\
Evolutionary example &
\textit{Trichoplax}-adjacent pre-bilaterian &
\textit{Dickinsonia}, \textit{Kimberella} \\
Developmental equivalent &
Blastula (radial) &
Early gastrula (bilateral axis established) \\
Independent prediction & Yes --- Stage~1 result
alone is publishable & Yes --- Stage~2 result
is the Dickinsonia result \\
\bottomrule
\end{tabular}
\end{center}

\subsection{Why Wnt alone is insufficient for
$\mathcal{A}_{E2}$}

The literature (2022--2024) confirms uniformly:

\begin{enumerate}
  \item Wnt/$\beta$-catenin activation in non-bilaterian
        organisms (\textit{Hydra}, planarians,
        \textit{Clytia}, \textit{Trichoplax}) produces
        ectopic primary axes, oral identity transformations,
        and axis elongation --- but not bilateral symmetry.

  \item Bilateral symmetry in all tested organisms
        requires the \textbf{orthogonal deployment of BMP}
        relative to the established Wnt axis. BMP
        specifies the dorsal-ventral axis independently
        of the AP axis.

  \item BMP pathway components (\textit{bmp2/4}, ADMP,
        Chordin/noggin as BMP antagonist establishing
        the gradient) are present in \textit{Trichoplax}
        \cite{srivastava2008} and all cnidarians but are
        not deployed orthogonally in non-bilaterian
        organisms under normal conditions.

  \item The $\mathcal{A}_{E1} \to \mathcal{A}_{E2}$
        transition in evolutionary history was the event
        in which BMP was first deployed orthogonal to the
        Wnt axis. This is the molecular event that opened
        the bilateral symmetry dimension.
\end{enumerate}

\begin{remark}
The original Dickinsonia Protocol (Section~8 of the
original pre-registration) specified:

\begin{quote}
\textit{``2. Inhibit further mesoderm specification
programmes (BMP gradient manipulation).''}
\end{quote}

This step was present but inverted: the original protocol
specified BMP \textbf{inhibition}, not activation. The
geometric derivation in this amendment corrects this.
BMP must be \textbf{activated orthogonal to the Wnt axis}
to specify Axis~2. BMP inhibition prevents the second axis,
maintaining the organism in $\mathcal{A}_{E1}$.

This correction is significant. The original protocol
would have accessed $\mathcal{A}_{E1}$ (pre-bilateral
Ediacaran disc morphology), not $\mathcal{A}_{E2}$
(\textit{Dickinsonia}-grade bilateral morphology).
Both are valid and publishable results. Only $\mathcal{A}_{E2}$
corresponds to \textit{Dickinsonia}.
\end{remark}

\subsection{The BMP activation geometry}

\begin{definition}[Orthogonal BMP Deployment]
Orthogonal BMP deployment is the specification of a
BMP activity gradient along an axis perpendicular to
the established Wnt/$\beta$-catenin axis. This requires:
\begin{enumerate}[label=(\roman*)]
  \item A pre-established Wnt axis (Step~1 of the
        revised protocol must precede Step~2).
  \item BMP ligand presentation or agonist delivery
        localised to one side of the organism relative
        to the established Wnt axis --- not distributed
        radially.
  \item The BMP antagonist Chordin/noggin to establish
        the DV gradient (high BMP one side, low BMP
        the other, mediated by Chordin).
\end{enumerate}
\end{definition}

This geometry is experimentally implementable in
\textit{Trichoplax} using:
\begin{itemize}
  \item Step~1: GSK-3$\beta$ inhibitor (e.g.\ CHIR99021,
        $\sim$\$45) to activate Wnt/$\beta$-catenin and
        establish Axis~1.
  \item Step~2 (after Axis~1 is established, $\sim$24h):
        Localised BMP2/4 protein (\$60--120 per experiment)
        applied asymmetrically relative to the established
        Wnt axis, combined with noggin on the opposite
        side to establish the DV gradient.
\end{itemize}

% ══════════════════════════════════════════════════════════════
\section{The Revised Dickinsonia Protocol}
% ══════════════════════════════════════════════════════════════

This section replaces Section~8
(\textit{The Dickinsonia Protocol}) of the original
pre-registration in its entirety for the purposes of the
experimental execution. The original section stands in the
permanent record. This section is the operative protocol.

\subsection{Target organisms: two stages}

The experiment now targets two independently falsifiable
attractor states:

\textbf{Stage 1 target ($\mathcal{A}_{E1}$):}
Pre-bilateral Ediacaran multicellular form. Disc or sheet
morphology with a single specified oral-aboral or AP axis.
No bilateral symmetry. 2--4 cell types. Radial or
near-radial organisation. Consistent with Ediacaran disc
organisms (pre-\textit{Dickinsonia} grade).

\textbf{Stage 2 target ($\mathcal{A}_{E2}$):}
Bilateral Ediacaran multicellular form. Two orthogonal
body axes. Distinct anterior-posterior AND
dorsal-ventral polarity. 2--6 cell types. No organ systems.
Modular accretion growth along the primary axis. Consistent
with \textit{Dickinsonia costata} and
\textit{Kimberella} ($\sim$558--555\,Ma).

\subsection{Recommended starting material}

\textit{Trichoplax adhaerens}.

Rationale: \textit{Trichoplax} is the only living animal
confirmed to occupy an attractor immediately adjacent to
$\mathcal{A}_{E1}$. It already has oral-aboral polarity.
It already has Wnt pathway components. It already has BMP
pathway components. The molecular toolkit for both steps
is present and confirmed by genome sequencing
\cite{srivastava2008}. No genetic modification is required.

Alternative starting material for cross-species
convergence: any cnidarian polyp (e.g.\ \textit{Hydra
vulgaris}) where Wnt-axis specification is confirmed.

\subsection{Revised protocol: two-step}

\textbf{Environment (both steps):}
\begin{itemize}[noitemsep]
  \item Seawater-based medium, salinity 30--35 ppt
  \item pH 7.8--8.1
  \item Temperature: 20--22$^\circ$C
  \item Oxygen: $<$1\% O$_2$ (microaerophilic; Ediacaran
        shallow marine approximation)
  \item No light stimulus during arrest phase
\end{itemize}

\vspace{0.5em}
\textbf{Step 1 --- Axis 1 specification (Wnt activation):}

\begin{enumerate}[label=\arabic*.]
  \item Culture \textit{Trichoplax} in standard
        \textit{Trichoplax} medium to confluence.
  \item Add GSK-3$\beta$ inhibitor CHIR99021 at
        5--10\,$\mu$M. This activates
        Wnt/$\beta$-catenin signalling by preventing
        $\beta$-catenin degradation.\\
        \textit{Alternative:} Exogenous recombinant
        Wnt3a protein at 100--300\,ng/mL.
  \item Incubate 24\,h at 20$^\circ$C.
  \item Confirm Axis~1 specification by immunostaining
        for $\beta$-catenin nuclear localisation at
        one pole of the organism relative to the
        other. This is the Axis~1 establishment
        confirmation.
  \item \textbf{Stage 1 result is now observable and
        should be photographed and recorded before
        proceeding to Step~2.} This is the
        $\mathcal{A}_{E1}$ result --- the pre-bilateral
        Ediacaran attractor.
\end{enumerate}

\vspace{0.5em}
\textbf{Step 2 --- Axis 2 specification (BMP activation,
orthogonal):}

\begin{enumerate}[label=\arabic*.]
  \item After Axis~1 confirmation, apply
        recombinant human BMP2 or BMP4
        (50--200\,ng/mL) localised to one edge of the
        organism \textbf{perpendicular} to the
        established $\beta$-catenin gradient.\\
        Practical method: apply BMP2/4 via
        protein-soaked bead (heparin acrylic bead,
        $\sim$\$30) placed at one lateral margin.
  \item Simultaneously apply Chordin protein or
        noggin (50--100\,ng/mL) at the opposite
        lateral margin to establish the DV gradient
        (high BMP one side, low BMP the other).
  \item Incubate further 48--72\,h at 20$^\circ$C,
        $<$1\% O$_2$.
  \item Monitor every 12\,h by light microscopy for:
        elongation along the Wnt axis, bilateral
        organisation perpendicular to the BMP
        gradient, emergence of distinct dorsal and
        ventral surfaces.
\end{enumerate}

\textbf{Revised estimated cost:} \$550--750 per full run
(both steps). Step~1 alone: $\sim$\$200--300.

\textbf{Timeline:}
\begin{itemize}[noitemsep]
  \item Step~1 complete: 24--36\,h from start
  \item Step~2 complete: 72--108\,h from start
  \item Total: 4--5 days to bilateral attractor result
\end{itemize}

\subsection{Validation criteria}

\textbf{Stage 1 ($\mathcal{A}_{E1}$) validation:}
\begin{itemize}[noitemsep]
  \item Multicellular aggregate maintained; no
        organ systems; 2--4 cell types morphologically
        distinguishable
  \item $\beta$-catenin nuclear localisation
        asymmetric: one pole enriched relative to the
        other, confirming Axis~1 polarity
  \item Disc or sheet morphology with single-axis
        polarity
  \item Organism viable at 36\,h
\end{itemize}

\textbf{Stage 2 ($\mathcal{A}_{E2}$) validation:}
\begin{itemize}[noitemsep]
  \item Elongated form along the Wnt axis
  \item Distinct dorsal (BMP-high) and ventral
        (BMP-low / Chordin-high) surfaces
        morphologically distinguishable
  \item Bilateral symmetry: left-right mirror
        symmetry about the AP axis
  \item 2--6 morphologically distinct cell types
  \item No organ systems; no nervous system; no
        circulatory system
  \item Modular segmental organisation along the
        AP axis (if visible at this scale)
  \item Organism viable at 96\,h
\end{itemize}

% ══════════════════════════════════════════════════════════════
\section{New Pre-Registered Predictions}
% ══════════════════════════════════════════════════════════════

The following predictions are additions to the original
pre-registration. They are pre-specified before any
experimental result. They supplement (do not replace) the
Dickinsonia predictions in Section~9 of the original.

\begin{prediction}[Stage~1: Pre-Bilateral Ediacaran
Attractor ($\mathcal{A}_{E1}$)]
\label{pred:e1}
Activation of Wnt/$\beta$-catenin signalling in
\textit{Trichoplax adhaerens} (Step~1 of the revised
protocol) will produce a stable multicellular form with a
single specified oral-aboral or anterior-posterior axis,
disc or sheet morphology, and no bilateral symmetry.

This is the pre-bilateral Ediacaran attractor
$\mathcal{A}_{E1}$.

\textbf{Measurable output:} Asymmetric $\beta$-catenin
localisation at one pole confirmed by immunostaining;
multicellular disc or sheet form maintained for
$\geq$36\,h; no bilateral axis; no bilateral elongation.

\textbf{Geometric basis:} Wnt activation specifies Axis~1
(AP/oral-aboral polarity) only. $\mathcal{A}_{E1}$ is
the single-axis attractor. BMP has not been deployed;
Axis~2 does not exist. The organism is at the
pre-bilateral Ediacaran attractor by definition.

\textbf{Significance:} $\mathcal{A}_{E1}$ is an
independently falsifiable Ediacaran-grade attractor
intermediate. Its confirmation validates the two-axis
attractor model before Step~2 is run. A Stage~1
confirmation is a publishable result independently of
whether Stage~2 succeeds.

\textbf{Failure condition:} Wnt activation produces no
organisational change, or immediately produces bilateral
symmetry without BMP deployment, or produces organ system
complexity inconsistent with the Ediacaran grade.
\end{prediction}

\begin{prediction}[Stage~2: Bilateral Ediacaran Attractor
($\mathcal{A}_{E2}$) ---
\textit{Dickinsonia}-grade]
\label{pred:e2}
Activation of Wnt/$\beta$-catenin signalling (Step~1)
followed by orthogonal BMP activation (Step~2) in
\textit{Trichoplax adhaerens} will produce a stable
multicellular form with bilateral symmetry: a defined
anterior-posterior axis, a defined orthogonal
dorsal-ventral axis, 2--6 morphologically distinct cell
types, no organ systems, and modular growth along the
primary axis.

This is the bilateral Ediacaran attractor $\mathcal{A}_{E2}$
--- the \textit{Dickinsonia}-grade organism.

\textbf{Measurable output:} Elongated bilateral form;
morphologically distinct dorsal and ventral surfaces;
left-right mirror symmetry about the AP axis; 2--6
cell types; no organ systems; organism viable at 96\,h.

\textbf{Geometric basis:} $\mathcal{A}_{E2}$ is the
two-axis attractor. It requires both Axis~1 (Wnt) and
Axis~2 (BMP, orthogonal). Both axes are specified
sequentially in the protocol. By the Reproduction Theorem,
this is the Ediacaran bilateral attractor that every
animal developmental program passes through at the
blastula-to-gastrula transition. It is not lost.
It is accessible.

\textbf{Failure condition:} Step~2 does not produce
bilateral organisation after confirmed Step~1 Axis~1
specification; or produces organ system complexity
inconsistent with the Ediacaran grade; or produces
morphology not bilateral despite orthogonal BMP
deployment with confirmed gradient.
\end{prediction}

\begin{prediction}[Two-Axis Attractor Convergence from
Mechanistically Independent Starting Points]
\label{pred:convergence}
The $\mathcal{A}_{E2}$ bilateral Ediacaran attractor,
accessed via the two-step Wnt+BMP protocol in
\textit{Trichoplax}, will be morphologically
indistinguishable from the bilateral form produced by
the same two-step protocol in a cnidarian starting
organism (e.g.\ \textit{Hydra vulgaris}).

\textbf{Geometric basis:} $\mathcal{A}_{E2}$ is defined
by the geometry of two orthogonal axes, not by the
molecular identity of the starting organism. Different
starting organisms with different lineage-specific gene
content will converge on the same two-axis attractor
geometry, because the attractor is a geometric property
of the body plan space, not of the genome.

This is the Ediacaran-grade analogue of the Cross-Kingdom
Convergence Prediction (Prediction 2 in the original
pre-registration).

\textbf{Failure condition:} \textit{Trichoplax} and
\textit{Hydra} two-step protocols produce systematically
different bilateral morphologies distinguishable above
chance by blinded observers.
\end{prediction}

\begin{prediction}[BMP Inhibition Locks the Organism
in $\mathcal{A}_{E1}$]
\label{pred:lock}
After confirmed Step~1 Axis~1 specification in
\textit{Trichoplax}, addition of a BMP inhibitor
(LDN-193189, $\sim$\$40; or dorsomorphin) instead of
Step~2 BMP activation will maintain the organism in
$\mathcal{A}_{E1}$ (pre-bilateral disc morphology)
and prevent bilateral organisation indefinitely.

\textbf{Geometric basis:} If the $\mathcal{A}_{E1}
\to \mathcal{A}_{E2}$ transition requires BMP
deployment, then BMP inhibition after Axis~1
specification must prevent the transition. This is
the negative control prediction for the two-axis
attractor model.

\textbf{Significance:} This prediction is the
experimental demonstration that $\mathcal{A}_{E1}$ is
a stable attractor, not a transient state. Organisms
can be maintained in $\mathcal{A}_{E1}$ indefinitely
by blocking the dimensional expansion.

\textbf{Failure condition:} BMP inhibition does not
prevent bilateral organisation after Axis~1 is
specified --- bilateral morphology emerges without
BMP regardless.
\end{prediction}

% ══════════════════════════════════════════════════════════════
\section{Connection to the LUCA Paper}
% ══════════════════════════════════════════���═══════════════════

This amendment is geometrically connected to the companion
publication \textit{The Ribosome Is the LUCA}
(DOI: \href{https://doi.org/10.5281/zenodo.18988844}
{10.5281/zenodo.18988844}).

The LUCA paper establishes that the FECA$\to$LECA
mitochondrial event is a dimensional expansion of the
cellular attractor space: the energy dimension opened,
and all post-LECA organisms are permanently in the
expanded energy attractor space.

Theorem~\ref{thm:bilateral} of this amendment establishes
that the $\mathcal{A}_{E1} \to \mathcal{A}_{E2}$
bilateral symmetry transition is also a dimensional
expansion, structurally identical to the FECA$\to$LECA
event but in morphological attractor space.

\begin{corollary}[Dimensional Expansions Are the
Structural Events of Macroevolution]
The geometry of the Attractor Framework identifies two
classes of evolutionary event:
\begin{enumerate}[label=(\roman*)]
  \item \textbf{Attractor traversal:} the organism moves
        through the attractor landscape along a known
        trajectory. The organism changes but remains in
        the same dimensional space. Accessible by
        developmental arrest.
  \item \textbf{Dimensional expansion:} a new independent
        dimension is added to the attractor space. The
        organism is now in a fundamentally larger space.
        Not reversible by arrest alone --- the new
        dimension must be actively closed.
\end{enumerate}

The known dimensional expansions in the history of life
accessible to this framework are:
\begin{itemize}[noitemsep]
  \item LUCA $\to$ LECA: energy dimension (mitochondria)
  \item $\mathcal{A}_{E1} \to \mathcal{A}_{E2}$:
        bilateral axis dimension (BMP orthogonal to Wnt)
\end{itemize}

All other major transitions in the evolutionary table
are attractor traversals, not dimensional expansions.
They are accessible by developmental arrest. Dimensional
expansion events are accessible by the two-step
protocol: establish the prior attractor, then activate
the new dimension.
\end{corollary}

% ══════════════════════════════════════════════════════════════
\section{Updated Falsifiability Conditions}
% ══════════════════════════════════════════════════════════════

The following conditions are added to the falsifiability
table of the original pre-registration (Table~4,
conditions F1--F7). They are additions, not replacements.
Conditions F1--F7 remain in force.

\begin{center}
\begin{tabularx}{\textwidth}{lX}
\toprule
\textbf{ID} & \textbf{Condition} \\
\midrule
F8 & Wnt activation in \textit{Trichoplax} produces
bilateral symmetry without BMP deployment. Would falsify
the two-axis attractor model and require the original
single-step Dickinsonia Protocol to be reinstated.
\\[0.4em]
F9 & Orthogonal BMP deployment after confirmed Axis~1
specification in \textit{Trichoplax} does not produce
bilateral organisation. Would falsify Prediction A1-P2
and require identification of an alternative second-axis
mechanism. Would not falsify the framework or the
Reproduction Theorem.
\\[0.4em]
F10 & BMP pathway components are absent or non-functional
in \textit{Trichoplax adhaerens} despite genome annotation
showing their presence. Would require switching to a
cnidarian starting organism and would not falsify the
framework.
\\[0.4em]
F11 & $\mathcal{A}_{E1}$ (disc morphology after Wnt
activation) is not a stable attractor --- the organism
immediately transits to $\mathcal{A}_{E2}$ upon Wnt
activation, or degrades, with no stable intermediate.
Would require revision of the two-stage model. Would not
falsify the framework or the Reproduction Theorem.
\\
\bottomrule
\end{tabularx}
\end{center}

\noindent None of conditions F8--F11 has been met as of
amendment date 2026-03-12. None can have been met: no
experiment has been run.

% ══════════════════════════════════════════════════════════════
\section{Summary: What Changed, What Did Not}
% ══════════════════════════════════════════════════════════════

\begin{center}
\begin{tabularx}{\textwidth}{lXX}
\toprule
\textbf{Element} &
\textbf{Status in original pre-registration} &
\textbf{Status after Amendment A1} \\
\midrule
Reproduction Theorem
  & Unchanged & Unchanged \\
Inverse Accessibility Theorem
  & Unchanged & Unchanged \\
LECA Resurrection Protocol
  & Unchanged & Unchanged \\
LECA predictions (P1--P5)
  & Unchanged & Unchanged \\
Cross-Kingdom Convergence Prediction
  & Unchanged & Unchanged \\
LINE-1 independent access prediction
  & Unchanged & Unchanged \\
LUCA geometric status
  & Unchanged & Unchanged \\
Falsifiability conditions F1--F7
  & Unchanged & Unchanged \\
\midrule
Dickinsonia Protocol (Section~8)
  & Single-step (Wnt only)
  & \textbf{Revised: two-step (Wnt + orthogonal BMP)} \\
Ediacaran attractor structure
  & Single target ($\mathcal{A}_{E2}$)
  & \textbf{Two stages: $\mathcal{A}_{E1}$ intermediate
    + $\mathcal{A}_{E2}$ bilateral target} \\
New predictions
  & None
  & \textbf{A1-P1 through A1-P4 added} \\
New falsifiability conditions
  & None
  & \textbf{F8--F11 added} \\
\bottomrule
\end{tabularx}
\end{center}

% ══════════════════════════════════════════════════════════════
\section{Locked Statement}
% ═════════════════════════��════════════════════════════════════

\begin{center}
\colorbox{warnbox}{
\begin{minipage}{0.88\textwidth}
\vspace{0.6em}
\textbf{This amendment is locked as of 2026-03-12.}

\medskip
The original pre-registration stands permanently
unchanged at DOI:
\href{https://doi.org/10.5281/zenodo.18986790}
{10.5281/zenodo.18986790}.

\medskip
This amendment refines one prediction in that record:
the Dickinsonia Protocol. It does not modify the
framework, any theorem, any LECA prediction, or any
other element of the original.

\medskip
\textbf{The refinement was derived geometrically.}
The Wnt/BMP literature was integrated into the attractor
framework and produced the two-axis attractor theorem.
The theorem produced the revised protocol. The revised
protocol produced four new pre-specified, falsifiable
predictions. All of this was completed before any
experiment was run.

\medskip
\textbf{The geometry required bilateral symmetry to have
two axes. The literature confirmed that it does. The
protocol was revised to match. This is the framework
working correctly.}

\medskip
\textit{``Bilateral symmetry is not a destination. It is
a second opening. The first axis is the ground.
The second axis is the expansion.
The geometry always required both.
Now the protocol states both explicitly.''}

\vspace{0.4em}
\hfill --- Eric Robert Lawson, March 12, 2026
\vspace{0.4em}
\end{minipage}
}
\end{center}

% ══════════════════════════════════════════════════════════════
\section*{Repository and Reproducibility}
\addcontentsline{toc}{section}{Repository and
Reproducibility}
% ══════════════════════════════════════════════════════════════

All derivation records, protocols, reasoning artifacts,
pre-registration documents, and this amendment are
publicly available at:

\begin{center}
\href{https://github.com/Eric-Robert-Lawson/attractor-oncology}
{\texttt{https://github.com/Eric-Robert-Lawson/attractor-oncology}}
\\[0.3em]
\small Repository ID: 1172048282 \textbar{}
MIT License \textbar{} Python \textbar{} Public
\end{center}

% ══════════════════════════════════════════════════════════════
\section*{Competing Interests}
\addcontentsline{toc}{section}{Competing Interests}
% ══════════════════════════════════════════════════════════════

The author has no competing interests to declare. No
institutional funding was received. No conflicts of
interest exist. All data used is publicly available.
All code is open source under the MIT licence.

% ── BIBLIOGRAPHY ──────────────────────────────────────────────
\newpage
\bibliographystyle{plainnat}
\begin{thebibliography}{99}

\bibitem{holzem2024}
Holzem, M., Boutros, M., \& Holstein, T.W. (2024).
The origin and evolution of Wnt signalling.
\textit{Nature Reviews Genetics}, 25(7), 500--512.
\doi{10.1038/s41576-024-00699-w}

\bibitem{clark2023}
Clark, E.M., \& Petersen, C.P. (2023).
BMP suppresses WNT to integrate patterning of
orthogonal body axes in adult planarians.
\textit{PLOS Genetics}, 19(2), e1010608.
\doi{10.1371/journal.pgen.1010608}

\bibitem{morsdorf2024}
M\"{o}rsdorf, D., et al. (2024).
Highly conserved and extremely evolvable: BMP
signalling in secondary axis patterning of
Cnidaria and Bilateria.
\textit{Discover Developmental Biology}, 1, 7.
\doi{10.1007/s00427-024-00714-4}

\bibitem{uveira2025}
Uveira, J., et al. (2025).
Planar cell polarity coordination in a cnidarian
embryo provides clues to animal body axis evolution.
\textit{eLife}, e104508.
\doi{10.7554/eLife.104508}

\bibitem{srivastava2008}
Srivastava, M., et al. (2008).
The \textit{Trichoplax} genome and the nature of
placozoans.
\textit{Nature}, 454(7207), 955--960.
\doi{10.1038/nature07191}

\bibitem{bobrovskiy2018}
Bobrovskiy, I., et al. (2018).
Ancient steroids establish the Ediacaran fossil
\textit{Dickinsonia} as one of the earliest animals.
\textit{Science}, 361(6408), 1246--1249.
\doi{10.1126/science.aat7228}

\bibitem{prereg}
Lawson, E.R. (2026).
Attractor Geometry of Life: Pre-Registration of
Geometric Derivations, Experimental Predictions,
and Falsifiability Conditions.
\textit{Zenodo}.
\doi{10.5281/zenodo.18986790}

\bibitem{luca2026}
Lawson, E.R. (2026).
The Ribosome Is the LUCA: A Geometric Derivation
That the Last Universal Common Ancestor Is Not
Extinct, Has Never Been Extinct, and Is Currently
Executing in Every Cell on Earth.
\textit{Zenodo}.
\doi{10.5281/zenodo.18988844}

\bibitem{moody2024}
Moody, E.R.R., et al. (2024).
Metabolism, genome and age of the last universal
common ancestor.
\textit{Nature}, 628, 545--550.
\doi{10.1038/s41559-024-02461-1}

\bibitem{kalinka2010}
Kalinka, A.T., et al. (2010).
Gene expression divergence recapitulates the
developmental hourglass model.
\textit{Nature}, 468(7325), 811--814.
\doi{10.1038/nature09634}

\bibitem{irie2011}
Irie, N., \& Kuratani, S. (2011).
Comparative transcriptome analysis reveals vertebrate
phylotypic period during organogenesis.
\textit{Nature Communications}, 2, 248.
\doi{10.1038/ncomms1248}

\end{thebibliography}

\end{document}
